% ENEM 2014-2016 — 11 questões MCQ — Porcentagem
% Fonte: lista "Porcentagem II" do site projetoagathaedu.com.br
% Omitidas: Q2,Q3,Q4,Q8,Q13,Q15 (imagens)

---
tipo: Múltipla Escolha
dificuldade: Média
nivel: médio
disciplina: Matemática
assunto: Porcentagem
gabarito: C
tags: [ENEM, porcentagem, gotas, soro, saúde, reidratação]
fonte: concurso
concurso: ENEM
banca: INEP
ano: 2016
numero: 1
---
\question
Um paciente necessita de reidratação endovenosa feita por meio de cinco frascos de soro durante 24 h. Cada frasco tem um volume de 800 mL de soro. Nas primeiras quatro horas, deverá receber 40\% do total a ser aplicado. Cada mililitro de soro corresponde a 12 gotas.

O número de gotas por minuto que o paciente deverá receber após as quatro primeiras horas será
\begin{choices}
\choice 16
\choice 20
\CorrectChoice 24
\choice 34
\choice 40
\end{choices}

---
tipo: Múltipla Escolha
dificuldade: Média
nivel: médio
disciplina: Matemática
assunto: Porcentagem
gabarito: E
tags: [ENEM, porcentagem, censo, renda, desigualdade]
fonte: concurso
concurso: ENEM
banca: INEP
ano: 2015
numero: 5
---
\question
Segundo dados apurados no Censo 2010, para uma população de 101,8 milhões de brasileiros com 10 anos ou mais de idade e que teve algum tipo de rendimento em 2010, a renda média mensal apurada foi de R\$ 1.202,00. A soma dos rendimentos mensais dos 10\% mais pobres correspondeu a apenas 1,1\% do total de rendimentos dessa população considerada, enquanto que a soma dos rendimentos mensais dos 10\% mais ricos correspondeu a 44,5\% desse total.

Qual foi a diferença, em reais, entre a renda média mensal de um brasileiro que estava na faixa dos 10\% mais ricos e de um brasileiro que estava na faixa dos 10\% mais pobres?
\begin{choices}
\choice 240,40
\choice 548,11
\choice 1.723,67
\choice 4.026,70
\CorrectChoice 5.216,68
\end{choices}

---
tipo: Múltipla Escolha
dificuldade: Difícil
nivel: médio
disciplina: Matemática
assunto: Porcentagem
gabarito: D
tags: [ENEM, porcentagem, juros, financiamento, prestações, saldo devedor]
fonte: concurso
concurso: ENEM
banca: INEP
ano: 2015
numero: 6
---
\question
Um casal realiza um financiamento imobiliário de R\$ 180\,000,00, a ser pago em 360 prestações mensais, com taxa de juros efetiva de 1\% ao mês. A primeira prestação é paga um mês após a liberação dos recursos e o valor da prestação mensal é de R\$ 500,00 mais juro de 1\% sobre o saldo devedor (valor devido antes do pagamento). Observe que, a cada pagamento, o saldo devedor se reduz em R\$ 500,00 e considere que não há prestação em atraso.

Efetuando o pagamento dessa forma, o valor, em reais, a ser pago ao banco na décima prestação é de
\begin{choices}
\choice 2\,075,00.
\choice 2\,093,00.
\choice 2\,138,00.
\CorrectChoice 2\,255,00.
\choice 2\,300,00.
\end{choices}

---
tipo: Múltipla Escolha
dificuldade: Fácil
nivel: médio
disciplina: Matemática
assunto: Porcentagem
gabarito: A
tags: [ENEM, porcentagem, fração, representação]
fonte: concurso
concurso: ENEM
banca: INEP
ano: 2015
numero: 7
---
\question
Uma pesquisa recente aponta que 8 em cada 10 homens brasileiros dizem cuidar de sua beleza, não apenas de sua higiene pessoal.

Outra maneira de representar esse resultado é exibindo o valor percentual dos homens brasileiros que dizem cuidar de sua beleza. Qual é o valor percentual que faz essa representação?
\begin{choices}
\CorrectChoice 80\%
\choice 8\%
\choice 0,8\%
\choice 0,08\%
\choice 0,008\%
\end{choices}

---
tipo: Múltipla Escolha
dificuldade: Média
nivel: médio
disciplina: Matemática
assunto: Porcentagem
gabarito: E
tags: [ENEM, porcentagem, potência, proporcionalidade, fluxo sanguíneo, Poiseuille]
fonte: concurso
concurso: ENEM
banca: INEP
ano: 2015
numero: 9
---
\question
O fisiologista francês Jean Poiseuille estabeleceu, na primeira metade do século XIX, que o fluxo de sangue por meio de um vaso sanguíneo em uma pessoa é diretamente proporcional à quarta potência da medida do raio desse vaso. Suponha que um médico, efetuando uma angioplastia, aumentou em 10\% o raio de um vaso sanguíneo de seu paciente.

O aumento percentual do fluxo por esse vaso está entre
\begin{choices}
\choice 7\% e 8\%
\choice 9\% e 11\%
\choice 20\% e 22\%
\choice 39\% e 41\%
\CorrectChoice 46\% e 47\%
\end{choices}

---
tipo: Múltipla Escolha
dificuldade: Fácil
nivel: médio
disciplina: Matemática
assunto: Porcentagem
gabarito: B
tags: [ENEM, porcentagem, saúde, nutrição, sódio, OMS]
fonte: concurso
concurso: ENEM
banca: INEP
ano: 2015
numero: 10
---
\question
A Organização Mundial da Saúde (OMS) recomenda que o consumo diário de sal de cozinha não exceda 5 g. Sabe-se que o sal de cozinha é composto por 40\% de sódio e 60\% de cloro.

Qual é a quantidade máxima de sódio proveniente do sal de cozinha, recomendada pela OMS, que uma pessoa pode ingerir por dia?
\begin{choices}
\choice 1\,250 mg
\CorrectChoice 2\,000 mg
\choice 3\,000 mg
\choice 5\,000 mg
\choice 12\,500 mg
\end{choices}

---
tipo: Múltipla Escolha
dificuldade: Média
nivel: médio
disciplina: Matemática
assunto: Porcentagem
gabarito: A
tags: [ENEM, porcentagem, variação percentual, receita, leite]
fonte: concurso
concurso: ENEM
banca: INEP
ano: 2015
numero: 11
---
\question
Um fornecedor vendia caixas de leite a um supermercado por R\$ 1,50 a unidade. O supermercado costumava comprar 3\,000 caixas de leite por mês desse fornecedor. Uma forte seca, ocorrida na região onde o leite é produzido, forçou o fornecedor a encarecer o preço de venda em 40\%. O supermercado decidiu então cortar em 20\% a compra mensal dessas caixas de leite. Após essas mudanças, o fornecedor verificou que sua receita nas vendas ao supermercado tinha aumentado.

O aumento da receita nas vendas do fornecedor, em reais, foi de
\begin{choices}
\CorrectChoice 540.
\choice 600.
\choice 900.
\choice 1\,260.
\choice 1\,500.
\end{choices}

---
tipo: Múltipla Escolha
dificuldade: Média
nivel: médio
disciplina: Matemática
assunto: Porcentagem
gabarito: A
tags: [ENEM, porcentagem, intervalo, composição, transparência, películas]
fonte: concurso
concurso: ENEM
banca: INEP
ano: 2014
numero: 12
---
\question
Os vidros para veículos produzidos por certo fabricante têm transparências entre 70\% e 90\%, dependendo do lote fabricado. Isso significa que, quando um feixe luminoso incide no vidro, uma parte entre 70\% e 90\% da luz consegue atravessá-lo. Os veículos equipados com vidros desse fabricante terão instaladas, nos vidros das portas, películas protetoras cuja transparência, dependendo do lote fabricado, estará entre 50\% e 70\%. Considere que uma porcentagem $P$ da intensidade da luz, proveniente de uma fonte externa, atravessa o vidro e a película.

De acordo com as informações, o intervalo das porcentagens que representam a variação total possível de $P$ é
\begin{choices}
\CorrectChoice $[35\,;\,63]$.
\choice $[40\,;\,63]$.
\choice $[50\,;\,70]$.
\choice $[50\,;\,90]$.
\choice $[70\,;\,90]$.
\end{choices}

---
tipo: Múltipla Escolha
dificuldade: Média
nivel: médio
disciplina: Matemática
assunto: Porcentagem
gabarito: B
tags: [ENEM, porcentagem, saneamento, esgoto, tratamento]
fonte: concurso
concurso: ENEM
banca: INEP
ano: 2014
numero: 14
---
\question
Uma organização não governamental divulgou um levantamento de dados realizado em algumas cidades brasileiras sobre saneamento básico. Os resultados indicam que somente 36\% do esgoto gerado nessas cidades é tratado, o que mostra que 8 bilhões de litros de esgoto sem nenhum tratamento são lançados todos os dias nas águas.

Uma campanha para melhorar o saneamento básico nessas cidades tem como meta a redução da quantidade de esgoto lançado nas águas diariamente, sem tratamento, para 4 bilhões de litros nos próximos meses.

Se o volume de esgoto gerado permanecer o mesmo e a meta dessa campanha se concretizar, o percentual de esgoto tratado passará a ser
\begin{choices}
\choice 72\%
\CorrectChoice 68\%
\choice 64\%
\choice 54\%
\choice 18\%
\end{choices}

---
tipo: Múltipla Escolha
dificuldade: Fácil
nivel: médio
disciplina: Matemática
assunto: Porcentagem
gabarito: D
tags: [ENEM, porcentagem, agricultura, território, Brasil]
fonte: concurso
concurso: ENEM
banca: INEP
ano: 2014
numero: 16
---
\question
O Brasil é um país com uma vantagem econômica clara no terreno dos recursos naturais, dispondo de uma das maiores áreas com vocação agrícola do mundo. Especialistas calculam que, dos 853 milhões de hectares do país, as cidades, as reservas indígenas e as áreas de preservação, incluindo florestas e mananciais, cubram por volta de 470 milhões de hectares. Aproximadamente 280 milhões se destinam à agropecuária, 200 milhões para pastagens e 80 milhões para a agricultura, somadas as lavouras anuais e as perenes, como o café e a fruticultura.

De acordo com os dados apresentados, o percentual correspondente à área utilizada para agricultura em relação à área do território brasileiro é mais próximo de
\begin{choices}
\choice 32,8\%
\choice 28,6\%
\choice 10,7\%
\CorrectChoice 9,4\%
\choice 8,0\%
\end{choices}

---
tipo: Múltipla Escolha
dificuldade: Fácil
nivel: médio
disciplina: Matemática
assunto: Porcentagem
gabarito: C
tags: [ENEM, porcentagem, distribuição de carga, ponte, engenharia]
fonte: concurso
concurso: ENEM
banca: INEP
ano: 2014
numero: 17
---
\question
Uma ponte precisa ser dimensionada de forma que possa ter três pontos de sustentação. Sabe-se que a carga máxima suportada pela ponte será de 12 t. O ponto de sustentação central receberá 60\% da carga da ponte, e o restante da carga será distribuído igualmente entre os outros dois pontos de sustentação.

No caso de carga máxima, as cargas recebidas pelos três pontos de sustentação serão, respectivamente,
\begin{choices}
\choice 1,8 t; 8,4 t; 1,8 t.
\choice 3,0 t; 6,0 t; 3,0 t.
\CorrectChoice 2,4 t; 7,2 t; 2,4 t.
\choice 3,6 t; 4,8 t; 3,6 t.
\choice 4,2 t; 3,6 t; 4,2 t.
\end{choices}
